\honest{Universal Law}{Eliezer Yudkowsky}{2007}

Lavoisier discovered that breathing and burning work on the same principle, which I call one of the most startling unifications in the history of science, joining the mundane realm of matter with the sacred realm of life. \nb{The link was older; see the notes on ``Universal Fire.''} Newton's was greater: it unified the planets, in the heavens once thought the home of the gods, with a falling apple, and ``Newton's discovery gave rise to the notion of a universal law,'' the same ``everywhere and everywhen, with literally zero exceptions.'' \nb{Descartes, writing between 1629 and 1633, already called the rules by which matter changes ``the `laws of nature'.'' Newton's gravitation gave the idea its most striking example.}

Such a law seems absurd to human minds. Everyday categories have exceptions: a tiger does not act like a buffalo, and most buffalo have four legs, but perhaps this one has three. Only in moral rules do we want a law binding everyone, as when someone takes more than a fair share of meat; yet even there, if a stronger tribe threatened to spear us all unless Bob got double meat this once, we would give it. So a rule with no exceptions seems like the product of fanatics.

That is the customary accusation against scientists. But the universe really is ``insanely rigid.'' ``As far as we know,'' conservation of momentum has never once been violated. Our models rarely fail, and last a generation or two, not centuries; but that does not make the universe whimsical, which would confuse the map with the territory. When a model fails, the new one is again exceptionless, and the universe was following it all along, as general relativity governed Mercury's orbit for decades, and the collapse of stars for billions of years, before anyone knew. ``It is only our model that was mistaken''; that is, ``or so our new model tells us.''

I have only 80 per cent confidence that the lightspeed limit will last. That does not mean I think it holds 80 per cent of the time. What I give 80 per cent is that it is ``absolutely inviolable throughout the entirety of space and time.''

The ancient Greeks did not discover science partly because they did not think you could generalize from experiments: they cared about ``normal'' phenomena, and thought a contrived setup would give a ``monstrous'' result with no bearing on how things really work. \nb{No source is given. Galen, a Greek physician of the second century, tied off nerves and ureters in living animals to learn how voice and kidneys work, and applied the results to humans.}

So humans dream before they learn better. To think like reality, the universe that dreamed before there were humans, here is ``the Tao'': ``Since the beginning / not one unusual thing / has ever happened.'' \nb{No attribution is given, and the lines could not be traced to any earlier source.}
